% 数列の超体系的解説：付録 深化版
% 漸化式を出発点に、差分、反復、微分方程式、数値解析、
% 線形代数、母関数、チェビシェフ多項式、漸近解析、組合せ論へ進む。

\chapter{付録}

\section*{この付録で見るもの}

本編では、漸化式を与えられたとき、その一般項を求めるための方法を学んできた。
付録では、同じ漸化式をもう少し遠くから眺める。

漸化式は、次の項を決める短い規則である。
けれども、その短い規則の中には、いくつもの数学的な問いが隠れている。

\begin{itemize}
  \item 項と項の間では、どのような変化が起きているのか。
  \item 同じ操作を何度も繰り返すと、値はどこへ行くのか。
  \item 一般項が分からなくても、増加や収束について何が言えるのか。
  \item 数列全体を一つの式や一つの行列にまとめると、何が見えるのか。
  \item ある漸化式が、三角関数や多項式、組合せ論とつながるのはなぜか。
\end{itemize}

この付録では、これらを別々の話として紹介しない。
一つの数列を調べているうちに、必要になった考え方が次の理論を呼び込む、という順番で進む。

例えば、隣り合う項の差を調べると、差分という操作が現れる。
差分を細かくしたくなると微分が現れ、微分方程式を計算するためにオイラー法が現れる。
解を直接求めることが難しければ、解へ近づく数列を作りたくなり、ニュートン法が現れる。
数列全体を一つにまとめたくなれば母関数になり、複数の項をまとめたくなれば行列になる。

したがって、ここでの目標は用語を増やすことではない。
「こうしたいから、この考え方を導入する」という数学の動機を追いかけることである。

この付録では、話の流れに合わせて添字を $0$ から始めることがある。
本文と添字の約束が異なる場合は、その場で明示する。

\input{draft/appendix_pattern_cases_draft.tex}

\section{差分から始まる離散的な微積分}

\subsection{数列の変化を測る}

数列
\[
  a_0,a_1,a_2,\ldots
\]
の値だけを眺めていると、変化の規則が見えにくいことがある。
そこで、隣り合う二項の差を考える。
\[
  \Delta a_n=a_{n+1}-a_n
\]
これを数列の一階差分という。

例えば、
\[
  a_n=n^2
\]
ならば、
\[
  \Delta a_n=(n+1)^2-n^2=2n+1
\]
である。
さらに差分を取ると、
\[
  \Delta^2a_n
  =\Delta(2n+1)
  =(2(n+1)+1)-(2n+1)=2
\]
となる。

二次式が二回の差分で定数になった。
これは、二次関数を二回微分すると定数になることに対応している。
差分は、離散的な世界で変化の速さを測る操作なのである。

\subsection{差分を取ると次数が下がる理由}

多項式
\[
  p(n)=c_mn^m+c_{m-1}n^{m-1}+\cdots+c_0
\]
を考える。
二項定理から、
\[
  (n+1)^m-n^m
  =mn^{m-1}+\text{それより低い次数の項}
\]
である。
したがって、$\Delta p(n)$ の次数は、$p(n)$ の次数より一つ低い。

特に、$m$ 次多項式に $m$ 回差分を取ると定数になり、$m+1$ 回差分を取ると $0$ になる。
この事実から、数列が多項式で表されるかどうかを、差分表から推測できる。

\[
\begin{array}{c|cccccc}
n&0&1&2&3&4&5\\ \hline
a_n&0&1&4&9&16&25\\
\Delta a_n&1&3&5&7&9&\\
\Delta^2a_n&2&2&2&2&&
\end{array}
\]

この表は、値の列よりも、変化の列の方が単純になることを示している。
漸化式を変形するという本編の発想も、実は「単純な変化を探す」という作業だった。

\subsection{差分演算子とシフト演算子}

差分の仕組みを、記号で整理してみよう。
数列を一つの対象と考え、シフト演算子 $E$ を
\[
  (Ea)_n=a_{n+1}
\]
で定める。
すると、
\[
  \Delta a_n=a_{n+1}-a_n=(E-1)a_n
\]
であるから、
\[
  \Delta=E-1
\]
と書ける。

この記号の利点は、差分を何回も取る操作を展開できることである。
\[
  \Delta^2=(E-1)^2=E^2-2E+1
\]
だから、
\[
  \Delta^2a_n=a_{n+2}-2a_{n+1}+a_n
\]
となる。
さらに、
\[
  \Delta^k=(E-1)^k
  =\sum_{j=0}^{k}(-1)^{k-j}\binom{k}{j}E^j
\]
より、
\[
  \Delta^ka_n
  =\sum_{j=0}^{k}(-1)^{k-j}\binom{k}{j}a_{n+j}
\]
を得る。

ここで現れた二項係数は偶然ではない。
差分は「一つ先へ進める操作」と「何もしない操作」の差であり、何度も差を取ると二項定理が現れるのである。

\subsection{差分の逆操作としての和}

微分の逆操作が積分であるように、差分の逆操作は和である。
例えば、
\[
  \Delta a_n=b_n
\]
ならば、
\[
  a_{n+1}-a_n=b_n
\]
なので、$0$ から $n-1$ まで足し合わせると、
\[
  a_n-a_0=\sum_{k=0}^{n-1}b_k
\]
となる。
従って、
\[
  a_n=a_0+\sum_{k=0}^{n-1}b_k
\]
である。

これは階差型漸化式の解法そのものである。
本編では解法として使ったが、ここでは「差分を元に戻す操作」として理解できる。

例えば、
\[
  a_{n+1}-a_n=2n+1,\qquad a_0=0
\]
ならば、
\[
  a_n=\sum_{k=0}^{n-1}(2k+1)
\]
であり、これは
\[
  a_n=n^2
\]
を与える。
平方数は、奇数を順番に足してできるという見方も得られた。

\subsection{和の公式を差分で設計する}

逆に、和
\[
  1^2+2^2+\cdots+n^2
\]
を求めたいとする。
公式を覚えていなくても、差分を利用して公式の形を設計できる。

和を $S_n$ とおくと、
\[
  S_n-S_{n-1}=n^2
\]
である。
右辺が二次式だから、$S_n$ は三次式になると予想する。
\[
  S_n=An^3+Bn^2+Cn+D
\]
と置けば、
\[
\begin{aligned}
  S_n-S_{n-1}
  &=
  A\{n^3-(n-1)^3\}
  +B\{n^2-(n-1)^2\}
  +C\\
  &=
  3An^2+(-3A+2B)n+(A-B+C)
\end{aligned}
\]
である。
これが $n^2$ に一致するためには、
\[
  3A=1,\qquad -3A+2B=0,\qquad A-B+C=0
\]
である。
さらに $S_0=0$ より $D=0$ だから、
\[
  A=\frac13,\qquad B=\frac12,\qquad C=\frac16
\]
となる。
従って、
\[
  1^2+2^2+\cdots+n^2
  =\frac{n(n+1)(2n+1)}6
\]
が得られる。

ここで重要なのは、和の公式を暗記したことではない。
「差分を取ると元の被加数になるような式を探す」という、離散的な積分の発想を使ったことである。

\section{反復としての漸化式と力学系}

\subsection{数列を点の運動として見る}

一項前から次の項を決める漸化式
\[
  a_{n+1}=f(a_n)
\]
を考える。
初期値 $a_0$ を決めると、
\[
  a_0,\ f(a_0),\ f(f(a_0)),\ldots
\]
という列ができる。

この列を、関数 $f$ によって動かされる点の軌道と見る。
数学では、このように状態の更新規則を調べる分野を\textbf{力学系}という。
物理の運動を扱わなくても、値が次々に変化するなら力学系として考えられる。

この視点では、一般項よりも次の問いが中心になる。

\begin{itemize}
  \item 軌道はある値へ近づくか。
  \item ある範囲から飛び出すか。
  \item 周期的に繰り返すか。
  \item 初期値を少し変えると、軌道は大きく変わるか。
\end{itemize}

\subsection{不動点と周期点}

$f(L)=L$ を満たす点 $L$ を不動点という。
不動点から出発した軌道は、
\[
  L\longmapsto L\longmapsto L\longmapsto\cdots
\]
と動かない。

一方、$f^{\,p}(L)=L$ を満たし、$p$ 回より少ない反復では元に戻らない点を周期 $p$ の周期点という。
例えば、
\[
  f(x)=2-x
\]
では $1$ が不動点であり、$x\neq1$ ならば
\[
  x\longmapsto 2-x\longmapsto x
\]
となるので、周期 $2$ の軌道になる。

「不動点を探す」だけでは、数列の全体の動きは分からない。
不動点へ近づく軌道もあれば、そこを飛び越えて周期軌道になるものもある。

\subsection{一次関数による反復を完全に分類する}

最も基本的な反復として、
\[
  a_{n+1}=ra_n+b
\]
を考える。
$r\neq1$ のときの不動点は、
\[
  L=rL+b
\]
より、
\[
  L=\frac{b}{1-r}
\]
である。
そこで $u_n=a_n-L$ と置くと、
\[
\begin{aligned}
  u_{n+1}
  &=a_{n+1}-L\\
  &=ra_n+b-L\\
  &=r(a_n-L)\\
  &=ru_n
\end{aligned}
\]
となる。
つまり、不動点からのずれは等比数列になる。
\[
  a_n-L=r^n(a_0-L)
\]
である。

この一行から、動きが完全に分かる。

\begin{itemize}
  \item $|r|<1$ なら $a_n\to L$。
  \item $r>1$ なら、一般には不動点から遠ざかる。
  \item $r<-1$ なら、符号を交互に変えながら大きくなる。
  \item $r=-1$ なら、通常は周期 $2$ になる。
  \item $r=1$ なら、$a_{n+1}=a_n+b$ であり、等差数列になる。
\end{itemize}

ここでは、漸化式を解く方法と、反復の動きを分類する方法が一致した。
一般項を求めることが、動力学的な理解へ直結している例である。

\subsection{微分は不動点の安定性を測る}

一般の関数 $f$ について、不動点 $L$ の近くで $a_n=L+u_n$ と置く。
テイラー展開の一次近似から、
\[
  f(L+u_n)
  \approx f(L)+f'(L)u_n
  =L+f'(L)u_n
\]
となる。
したがって、
\[
  u_{n+1}\approx f'(L)u_n
\]
である。

不動点の近くでは、結局、一次関数の反復に近い。
そこで、
\[
  |f'(L)|<1
\]
ならばずれが縮み、$L$ は安定であると考えられる。
逆に、
\[
  |f'(L)|>1
\]
ならばずれが拡大し、$L$ は不安定になりやすい。

この判定は、一般項を求めなくても数列の極限を予測するための強力な道具になる。
ただし、これは $L$ の近くでの局所的な判定である。
初期値が遠い場合には、別の不動点へ行ったり、発散したりする可能性がある。

\subsection{ロジスティック写像：単純な規則から複雑な動きへ}

次の漸化式を考える。
\[
  a_{n+1}=r a_n(1-a_n)
\]
これは、人口の割合や資源の制限を簡単に表すモデルとして現れる。
$0<a_n<1$ なら、$a_n$ は割合として解釈できる。

右辺の $ra_n$ は、現在の人口に比例する増加を表す。
一方、$(1-a_n)$ は、人口が大きくなるほど増加の余地が小さくなることを表している。
一つの式の中に、増加と抑制が同時に入っている。

不動点を求めると、
\[
  L=rL(1-L)
\]
より、
\[
  L=0,\qquad L=1-\frac1r
\]
である。
後者における微分係数は、
\[
  f'(L)=r(1-2L)=2-r
\]
となる。
したがって、
\[
  |2-r|<1
\]
すなわち
\[
  1<r<3
\]
の範囲では、正の不動点が安定になると予想される。

ところが $r$ をさらに大きくすると、不動点へ収束せず、周期 $2$、周期 $4$、さらに複雑な軌道が現れる。
この現象を完全に説明するには、力学系の理論が必要になる。
しかし、ここで知っておきたいのは、非常に単純な一行の漸化式でも、一般項を求めるという問題を超えて、安定性や周期性という深い問題を生み出すということである。

\section{離散と連続を結ぶ}

\subsection{時間の幅を明示する}

微分方程式
\[
  \frac{dx}{dt}=F(x)
\]
を考える。
時間を幅 $h$ で区切り、
\[
  t_n=nh,\qquad a_n=x(t_n)
\]
と置く。
微分の定義から、
\[
  \frac{x(t_{n+1})-x(t_n)}h
  \approx \frac{dx}{dt}(t_n)
  =F(x(t_n))
\]
である。
従って、
\[
  a_{n+1}=a_n+hF(a_n)
\]
という漸化式が得られる。

$h$ は単なる記号ではない。
それは、一回の更新がどれくらいの時間を表しているかを決める。
もし $h=1$ と暗黙に置くなら、そのことを意識しておく必要がある。

\subsection{オイラー法の誤差を一歩だけ調べる}

テイラー展開により、
\[
  x(t+h)
  =x(t)+hx'(t)+\frac{h^2}{2}x''(t)+\cdots
\]
である。
オイラー法は、
\[
  x(t+h)\approx x(t)+hx'(t)
\]
と、二次以上の項を捨てている。
従って、一回の更新で生じる誤差は、通常 $h^2$ 程度である。

しかし、時間 $T$ まで進むには、およそ $T/h$ 回の更新が必要である。
一回あたり $h^2$ の誤差が累積すると、全体の誤差はおおよそ
\[
  \frac{T}{h}\cdot h^2=Th
\]
程度になる。
これが、オイラー法の全体誤差が $h$ の一次程度になる理由である。

厳密な誤差評価には条件が必要だが、細かい刻みを使うほど近似がよくなる理由はこのように説明できる。

\subsection{指数関数は連続な等比数列である}

最も基本的な微分方程式
\[
  \frac{dx}{dt}=\lambda x
\]
を考える。
オイラー法では、
\[
  a_{n+1}=a_n+h\lambda a_n
  =(1+h\lambda)a_n
\]
となる。
これは等比数列である。
\[
  a_n=(1+h\lambda)^n a_0
\]

一方、微分方程式の解は、
\[
  x(t)=e^{\lambda t}x(0)
\]
である。
$t=nh$ とすれば、
\[
  e^{\lambda nh}=\left(e^{\lambda h}\right)^n
\]
だから、連続的な指数関数も、細かい時間幅で見れば等比的な更新の極限として現れる。

実際、
\[
  e^{\lambda h}
  =1+\lambda h+\frac{\lambda^2h^2}{2}+\cdots
\]
なので、$h$ が小さいとき、
\[
  1+\lambda h
\]
は $e^{\lambda h}$ の一次近似である。

「等比数列」と「指数関数」は別々の公式ではない。
一段ずつ進む世界と、連続的に進む世界で、同じ増加の仕組みを記述している。

\subsection{離散的な安定性と刻み幅}

$\lambda<0$ のとき、連続解 $e^{\lambda t}$ は $0$ へ近づく。
オイラー法の漸化式
\[
  a_{n+1}=(1+h\lambda)a_n
\]
が $0$ へ近づくためには、
\[
  |1+h\lambda|<1
\]
が必要である。
従って、
\[
  0<h<-\frac2\lambda
\]
でなければならない。

元の微分方程式が安定でも、刻み幅を大きく取りすぎると、近似の漸化式は発散する。
これは、「連続化したから安心」ではないことを示している。
離散化は新しい漸化式を作る操作であり、その漸化式自身の性質を調べなければならない。

\section{解を直接求めず、近づく数列を設計する}

\subsection{テイラー展開は何をしているのか}

関数 $f(x)$ を $x=a$ の近くで調べる。
接線による近似は、
\[
  f(x)\approx f(a)+f'(a)(x-a)
\]
である。
二次の項まで含めれば、
\[
  f(x)
  \approx f(a)+f'(a)(x-a)
  +\frac{f''(a)}{2}(x-a)^2
\]
となる。

テイラー展開は、関数を多項式で近似する道具である。
ただし、「近似だから適当でよい」という意味ではない。
どの点の近くで、何次まで使うかを決めることで、関数の局所的な情報を数値計算に変換している。

\subsection{ニュートン法の導出}

方程式
\[
  f(x)=0
\]
の解を探す。
現在の近似値を $x_n$ とし、$x_n$ の近くでは関数を接線で置き換える。
\[
  0\approx f(x_n)+f'(x_n)(x-x_n)
\]
を $x$ について解くと、
\[
  x=x_n-\frac{f(x_n)}{f'(x_n)}
\]
である。
この値を次の近似値に採用する。
\[
  x_{n+1}=x_n-\frac{f(x_n)}{f'(x_n)}
\]

ニュートン法は、方程式を解く公式ではない。
方程式を、収束することを期待した漸化式へ変換する設計法である。

\subsection{平方根の例と二次収束}

\[
  f(x)=x^2-2
\]
にニュートン法を使うと、
\[
  x_{n+1}=\frac{x_n^2+2}{2x_n}
\]
となる。
$x_0=1$ とすれば、
\[
  x_1=\frac32,\quad
  x_2=\frac{17}{12},\quad
  x_3=\frac{577}{408}
\]
である。

$\alpha=\sqrt2$ とおくと、
\[
\begin{aligned}
  x_{n+1}-\alpha
  &=\frac{x_n^2+2}{2x_n}-\alpha\\
  &=\frac{x_n^2-2\alpha x_n+\alpha^2}{2x_n}\\
  &=\frac{(x_n-\alpha)^2}{2x_n}.
\end{aligned}
\]

誤差を $e_n=x_n-\alpha$ と書けば、
\[
  e_{n+1}=\frac{e_n^2}{2x_n}
\]
である。
誤差が二乗されるため、十分近くまで来ると正しい桁数が急速に増える。
この性質を二次収束という。

\subsection{一般の単純な根の近くで起こること}

$f(\alpha)=0$ かつ $f'(\alpha)\neq0$ とする。
$x_n=\alpha+e_n$ とおき、二次までテイラー展開すると、
\[
  f(x_n)=f'(\alpha)e_n+\frac{f''(\alpha)}2e_n^2+\cdots
\]
である。
同様に、
\[
  f'(x_n)=f'(\alpha)+f''(\alpha)e_n+\cdots
\]
となる。
これをニュートン法へ代入すると、一次の誤差が打ち消され、

\[
  e_{n+1}=C e_n^2+\text{三次以上の項}
\]
という形になる。
これがニュートン法が、根の近くで二次収束する理由である。

ただし、根が重なっていて $f'(\alpha)=0$ になる場合や、初期値が遠い場合には、この説明はそのまま使えない。
ニュートン法を適用する前に、何を近似しているのか、どの範囲で近似が有効なのかを確認する必要がある。

\subsection{微分を使わないニュートン法}

微分が計算しにくい場合、二つの近似値を使って接線の傾きを近似することができる。
\[
  f'(x_n)\approx
  \frac{f(x_n)-f(x_{n-1})}{x_n-x_{n-1}}
\]
と置くと、
\[
  x_{n+1}
  =
  x_n
  -
  f(x_n)
  \frac{x_n-x_{n-1}}{f(x_n)-f(x_{n-1})}
\]
を得る。
これは割線法と呼ばれる。

ニュートン法と割線法の違いは、単なる公式の違いではない。
ニュートン法は関数の傾きを正確に使い、割線法は過去の値から傾きを推測する。
どちらも「次の値をどう設計すれば、解へ近づくか」という同じ問題への答えである。

\section{線形漸化式と線形代数}

\subsection{解の集合には構造がある}

二階線形漸化式
\[
  a_{n+2}=pa_{n+1}+qa_n
\]
を考える。
二つの解 $a_n,b_n$ があれば、定数 $c,d$ に対して
\[
  c a_n+d b_n
\]
も解になる。
これを線形性という。

さらに、初期値 $a_0,a_1$ を指定すれば、以後の項は一意に決まる。
したがって、解の全体は「二つの初期値を選ぶ空間」と同じくらいの自由度を持つ。
このことを、解の空間の次元が $2$ である、と表現する。

特性方程式の二つの解から一般項を作る方法は、単に計算がうまくいくからではない。
二次の漸化式の解全体を、二つの基本的な解の組み合わせで表しているのである。

\subsection{特性方程式と固有値}

\[
  \boldsymbol v_n=
  \begin{pmatrix}a_{n+1}\\a_n\end{pmatrix}
\]
と置けば、
\[
  \boldsymbol v_{n+1}
  =
  \begin{pmatrix}p&q\\1&0\end{pmatrix}
  \boldsymbol v_n
\]
となる。
この行列を $A$ とする。

もし
\[
  A\boldsymbol v=\lambda\boldsymbol v
\]
を満たすベクトルがあれば、その方向の運動は
\[
  \boldsymbol v,\ \lambda\boldsymbol v,\ \lambda^2\boldsymbol v,\ldots
\]
となる。
行列による更新が、数の等比的な更新に変わった。

固有値を求める条件は、
\[
  \det(A-\lambda I)=0
\]
であり、
\[
  \lambda^2-p\lambda-q=0
\]
となる。
これは二階漸化式の特性方程式と一致する。
特性方程式は、行列の固有値を求める方程式としても現れているのである。

\subsection{重解のときに $n$ が現れる理由}

特性方程式が重解 $\lambda$ を持つ場合、二つの異なる等比数列だけでは解の自由度を埋められない。
例えば、
\[
  a_{n+2}-2a_{n+1}+a_n=0
\]
では、差分を取ると、
\[
  \Delta^2a_n=0
\]
である。
従って、
\[
  \Delta a_n=B
\]
は定数であり、
\[
  a_n=A+Bn
\]
となる。

特性方程式は
\[
  (\lambda-1)^2=0
\]
である。
重解 $1$ に対して、$1^n$ だけでなく $n1^n$ が現れた。
この $n$ は、差分を一回戻すときに現れる一次的な増加を記録している。

一般に、特性根 $r$ が $m$ 重に現れると、
\[
  r^n,\ nr^n,\ n^2r^n,\ldots,n^{m-1}r^n
\]
が解を作る。
大学で学ぶジョルダン標準形は、この現象を行列の側から整理したものである。

\subsection{複素数の根と振動}

特性根が
\[
  r=\rho(\cos\theta+i\sin\theta)
\]
と
\[
  \overline r=\rho(\cos\theta-i\sin\theta)
\]
であるとする。
オイラーの公式
\[
  e^{i\theta}=\cos\theta+i\sin\theta
\]
を使えば、二つの根は $\rho e^{i\theta}$ と $\rho e^{-i\theta}$ である。

このとき、実数解は
\[
  a_n=\rho^n\{C\cos(n\theta)+D\sin(n\theta)\}
\]
という形になる。
$\rho^n$ は振幅の増減を、$\cos(n\theta)$ と $\sin(n\theta)$ は振動を表す。

実根だけを考えていたときには見えなかった振動が、複素根の偏角として現れた。
複素数は、実数の数列に存在する振動を記録するための自然な言語になっている。

\section{チェビシェフ多項式：三角関数を多項式に翻訳する}

\subsection{同じ漸化式が二つの姿を持つ}

\[
  a_{n+1}=2x a_n-a_{n-1}
\]
を考える。
$x$ を固定すれば二階線形漸化式であり、特性方程式は
\[
  r^2-2xr+1=0
\]
である。

$-1\leq x\leq1$ のとき、$x=\cos\theta$ と書ける。
すると、
\[
  \cos((n+1)\theta)
  =
  2\cos\theta\cos(n\theta)-\cos((n-1)\theta)
\]
だから、$\cos(n\theta)$ は同じ漸化式を満たす。

ここで、「三角関数の値を計算する問題」と「漸化式を解く問題」が同じ式を持っている。
この一致を、変数 $x$ の多項式として保存したくなる。

\subsection{チェビシェフ多項式の定義と証明}

第一種チェビシェフ多項式 $T_n(x)$ を、
\[
  T_0(x)=1,\qquad T_1(x)=x
\]
および
\[
  T_{n+1}(x)=2xT_n(x)-T_{n-1}(x)
\]
で定める。

初めのいくつかは、
\[
\begin{aligned}
  T_0(x)&=1,\\
  T_1(x)&=x,\\
  T_2(x)&=2x^2-1,\\
  T_3(x)&=4x^3-3x,\\
  T_4(x)&=8x^4-8x^2+1
\end{aligned}
\]
である。

漸化式の右辺が多項式だけでできているので、すべての $T_n(x)$ は多項式になる。
一方、
\[
  T_n(\cos\theta)=\cos(n\theta)
\]
が成り立つ。

実際、$n=0,1$ では定義から明らかである。
$n$ と $n-1$ について成り立つと仮定すれば、
\[
\begin{aligned}
  T_{n+1}(\cos\theta)
  &=2\cos\theta T_n(\cos\theta)-T_{n-1}(\cos\theta)\\
  &=2\cos\theta\cos(n\theta)-\cos((n-1)\theta)\\
  &=\cos((n+1)\theta)
\end{aligned}
\]
となる。

三角関数を使って書かれた量を、$x$ の多項式に翻訳できた。
例えば、
\[
  \cos(3\theta)=4\cos^3\theta-3\cos\theta
\]
は、
\[
  T_3(x)=4x^3-3x
\]
という多項式の公式になる。

\subsection{三種類の領域：振動、境界、増大}

特性方程式
\[
  r^2-2xr+1=0
\]
の根は、
\[
  r=x\pm\sqrt{x^2-1}
\]
である。

\begin{itemize}
  \item $|x|<1$ なら根は複素数で、単位円上にある。$T_n(x)$ は振動する。
  \item $x=1$ または $x=-1$ では根が重なり、境界的な振る舞いになる。
  \item $|x|>1$ なら実根の絶対値が異なり、一方の根が支配的になる。
\end{itemize}

$x=\cos\theta$ の代わりに $x=\cosh t$ と置くと、
\[
  T_n(\cosh t)=\cosh(nt)
\]
となる。
三角関数の振動と双曲線関数の増大が、同じ多項式の異なる領域で現れている。

これは、特性根の絶対値が数列の増加を支配するという、線形漸化式の基本原理を別の言葉で見たものである。

\subsection{第二種チェビシェフ多項式}

第一種だけでなく、\[
  U_0(x)=1,\qquad U_1(x)=2x
\]
および
\[
  U_{n+1}(x)=2xU_n(x)-U_{n-1}(x)
\]
で定まる第二種チェビシェフ多項式も考えられる。
こちらは、
\[
  U_n(\cos\theta)
  =\frac{\sin((n+1)\theta)}{\sin\theta}
\]
を満たす。

第一種が余弦を記録するのに対して、第二種は正弦を記録する。
同じ漸化式でも初期値を変えると、異なる多項式族が現れる。
初期値は単なる付属条件ではなく、解の空間の中からどの解を選ぶかを決めている。

\subsection{直交性と近似への入口}

チェビシェフ多項式が特別なのは、三角関数を表すからだけではない。
余弦の直交性から、
\[
  \int_0^\pi\cos(m\theta)\cos(n\theta)\,d\theta=0
  \qquad(m\neq n)
\]
が得られる。

$x=\cos\theta$ と変数変換すれば、
\[
  \int_{-1}^{1}
  \frac{T_m(x)T_n(x)}{\sqrt{1-x^2}}\,dx=0
  \qquad(m\neq n)
\]
となる。
この性質を直交性という。

直交する多項式を使うと、複雑な関数を
\[
  c_0T_0(x)+c_1T_1(x)+c_2T_2(x)+\cdots
\]
のように展開して近似できる。
ここからフーリエ級数や数値計算、最良近似の理論へ進むことができる。

本書で必要なのは近似理論の完成ではない。
一つの漸化式が、三角関数を経由して、多項式の直交性や関数近似へつながるという道筋を見つけることである。

\clearpage
\section{具体的な型から特殊関数へ}

\subsection{「型を解く」ことの先にあるもの}

前半の実戦編では、式の形を見て変数を置き換えた。
定数項があれば不動点を引き、べきがあれば対数を取り、
三角関数型なら倍角公式を使った。

ここで、もう一歩だけ先へ進む。
その型が一つの問題に現れただけなら、解法を知っていれば十分である。
しかし、同じ型が波、確率、近似、幾何の中で何度も現れるなら、
その解法の中心にある関数を独立した数学的対象として研究したくなる。

\[
\begin{array}{c}
\text{具体的な漸化式}\\
\downarrow\quad\text{置換・公式・帰納法}\\
\text{標準的な解}\\
\downarrow\quad\text{別の問題にも繰り返し現れる}\\
\text{特殊関数}
\end{array}
\]

この流れを、いくつかの例で追う。

\subsection{階乗型の漸化式を実数へ延長する}

階乗は
\[
  (n+1)!=(n+1)n!
\]
を満たす。
この式の $n$ を実数 $x$ に広げたいなら、
\[
  F(x+1)=xF(x),\qquad F(1)=1
\]
を満たす関数を探すことになる。

この条件だけでは関数は一意に決まらない。
例えば、適切な周期関数を掛けても漸化式を保つからである。
そこで、正値性や滑らかさ、成長の仕方などの自然な条件を追加する。

代表的な選択が、
\[
  \Gamma(x)=\int_0^\infty t^{x-1}e^{-t}\,dt
\]
である。
部分積分により、
\[
\begin{aligned}
  \Gamma(x+1)
  &=\int_0^\infty t^xe^{-t}\,dt\\
  &=x\int_0^\infty t^{x-1}e^{-t}\,dt\\
  &=x\Gamma(x)
\end{aligned}
\]
となる。
さらに $\Gamma(1)=1$ なので、
\[
  \Gamma(n+1)=n!
\]
である。

ここでは、実戦編の「階比型で積を整理する」という発想が、
整数列の外へ延長されている。
階乗に似た数列を解くための道具だった積分表示が、
ガンマ関数という特殊関数になった。

\subsection{三角関数型からチェビシェフ多項式へ}

実戦編で扱った
\[
  a_{n+1}=2\cos\theta\,a_n-a_{n-1}
\]
に対して、
\[
  a_n=\cos(n\theta)
\]
を倍角公式と帰納法で示した。

ここで $\theta$ を毎回書く代わりに、
\[
  x=\cos\theta
\]
と置き、
\[
  T_n(x)=\cos(n\theta)
\]
と定義する。
すると、
\[
  T_{n+1}(x)=2xT_n(x)-T_{n-1}(x)
\]
となる。

つまり、三角関数型の漸化式を、$x$ の多項式の漸化式として読み直した。
これがチェビシェフ多項式である。

この変換が有効なのは、三角関数が多項式近似の問題に姿を変えるからである。
例えば、$[-1,1]$ 上で関数を近似したいとき、
\[
  1,x,x^2,\ldots
\]
をそのまま使う代わりに、
\[
  T_0(x),T_1(x),T_2(x),\ldots
\]
を使うと、区間上の振動を制御しやすい。

\subsection{同じ三項漸化式から、別の多項式族が生まれる}

チェビシェフ多項式だけが特別なのではない。
係数の異なる三項漸化式から、別の自然な多項式族が生まれる。

例えばルジャンドル多項式は、
\[
  (n+1)P_{n+1}(x)
  =(2n+1)xP_n(x)-nP_{n-1}(x)
\]
を満たす。
母関数は
\[
  \frac1{\sqrt{1-2xt+t^2}}
  =\sum_{n=0}^{\infty}P_n(x)t^n
\]
である。

この多項式は、
\[
  (1-x^2)y''-2xy'+n(n+1)y=0
\]
の多項式解でもある。
さらに、
\[
  \int_{-1}^{1}P_m(x)P_n(x)\,dx=0
  \qquad(m\neq n)
\]
という直交性を持つ。

一つの漸化式を解くという入口から、
\[
  \text{母関数}
  \longleftrightarrow
  \text{微分方程式}
  \longleftrightarrow
  \text{直交性}
\]
という三つの見方が得られる。

\subsection{ガウス関数の微分からエルミート多項式へ}

正規分布に現れる
\[
  e^{-x^2}
\]
を何度も微分すると、
必ず多項式と $e^{-x^2}$ の積になる。
そこで、
\[
  H_n(x)e^{-x^2}
  =(-1)^n\frac{d^n}{dx^n}e^{-x^2}
\]
で $H_n$ を定める。

生成関数は
\[
  e^{2xt-t^2}
  =\sum_{n=0}^{\infty}H_n(x)\frac{t^n}{n!}
\]
であり、$t$ で微分して係数を比較すると、
\[
  H_{n+1}(x)=2xH_n(x)-2nH_{n-1}(x)
\]
が得られる。

ここでは「母関数を作る」という発展パターンが、
直交多項式の具体的な漸化式を作り出している。
ルジャンドル多項式では区間 $[-1,1]$、
エルミート多項式では全実数と重み $e^{-x^2}$ が現れる。
問題の対称性や重みが変わると、自然な特殊関数も変わる。

\subsection{円筒の波からベッセル関数へ}

円筒座標で波動方程式を解くと、
\[
  x^2y''+xy'+(x^2-\nu^2)y=0
\]
が現れる。
その解であるベッセル関数は、整数次数について
\[
  J_{n+1}(x)
  =\frac{2n}{x}J_n(x)-J_{n-1}(x)
\]
を満たす。

三角関数型や直交多項式型と同じく、
隣り合う次数の関数を結ぶ漸化式がある。
ただし、ここでは解は多項式ではなく、円筒の波を表す関数になる。

生成関数
\[
  \exp\left(\frac{x}{2}\left(t-\frac1t\right)\right)
  =\sum_{n=-\infty}^{\infty}J_n(x)t^n
\]
を $t$ で微分して係数比較を行うと、
\[
  \frac{x}{2}\{J_{n-1}(x)+J_{n+1}(x)\}
  =nJ_n(x)
\]
が得られ、上の漸化式が戻る。

ここで、実戦編の「隣の項を結ぶ漸化式を利用する」という発想が、
波動方程式の解を効率よく計算する道具になっている。

\subsection{具体的な型と特殊関数の対応}

\begin{center}
\begin{tabular}{c|c|c}
  実戦での見抜き方 & 発展した対象 & 広がる理論\\ \hline
  階比型・積の整理 & ガンマ関数 & 積分・確率\\
  三角関数型 & チェビシェフ多項式 & 近似・直交性\\
  三項漸化式 & ルジャンドル多項式 & 球面対称性\\
  ガウス関数の微分 & エルミート多項式 & 正規分布・量子力学\\
  円筒の波の漸化式 & ベッセル関数 & 波動・境界値問題\\
\end{tabular}
\end{center}

この表で重要なのは、関数の名前を増やすことではない。
目の前の漸化式に対して行った
\[
  \text{置く}\quad
  \text{変形する}\quad
  \text{生成する}\quad
  \text{直交化する}
\]
という操作が、別の問題でも繰り返し現れた結果として、
特殊関数が生まれたということである。

\section{母関数：無限個の数を一つの対象にする}

\subsection{なぜ数列を式に詰め込むのか}

数列は無限個の項を持つ。
項を一つずつ扱う代わりに、すべての項を係数として一つの式にまとめる。
\[
  A(x)=a_0+a_1x+a_2x^2+\cdots
\]
これを母関数という。

母関数では、$x^n$ の係数が $a_n$ である。
したがって、母関数を変形して係数を読むことができれば、数列の問題を式の問題へ移せる。

\subsection{形式的冪級数としての母関数}

ここで、母関数を通常の関数と考える必要はない。
例えば、
\[
  1+x+x^2+x^3+\cdots
\]
は、$|x|<1$ では $1/(1-x)$ に収束するが、形式的冪級数としては、係数を並べた無限列である。

形式的に、
\[
  (1-x)(1+x+x^2+\cdots)=1
\]
と係数を比較できる。
収束するかどうかを考えなくても、漸化式の係数操作が可能になる。

一方、実際に数値を代入して関数として使う場合には、収束範囲を確認する必要がある。
同じ記号でも、形式的な議論と解析的な議論を区別しなければならない。

\subsection{線形漸化式を母関数へ移す}

\[
  a_{n+2}=a_{n+1}+a_n
\]
とし、
\[
  A(x)=\sum_{n=0}^{\infty}a_nx^n
\]
と置く。
漸化式に $x^n$ を掛けて足すと、
\[
  \sum_{n=0}^{\infty}a_{n+2}x^n
  =
  \sum_{n=0}^{\infty}a_{n+1}x^n
  +
  \sum_{n=0}^{\infty}a_nx^n
\]
である。

左辺と右辺を $A(x)$ で表すと、
\[
  \frac{A(x)-a_0-a_1x}{x^2}
  =
  \frac{A(x)-a_0}{x}+A(x)
\]
となる。
これを整理して、
\[
  A(x)=\frac{a_0+(a_1-a_0)x}{1-x-x^2}
\]
を得る。

分母の多項式 $1-x-x^2$ は、漸化式の特性方程式と同じ情報を持つ。
母関数は、特性方程式を別の形で記録しているのである。

\subsection{畳み込みが積になる}

母関数の特に重要な性質は、積の係数である。
\[
  A(x)=\sum_{n=0}^{\infty}a_nx^n,\qquad
  B(x)=\sum_{n=0}^{\infty}b_nx^n
\]
とすると、
\[
  A(x)B(x)
  =
  \sum_{n=0}^{\infty}
  \left(\sum_{k=0}^{n}a_kb_{n-k}\right)x^n
\]
となる。

内側の和
\[
  \sum_{k=0}^{n}a_kb_{n-k}
\]
を畳み込みという。
二つの選択を組み合わせる数え上げでは、畳み込みが自然に現れる。
母関数は、畳み込みを積へ変えることで、組合せ論の漸化式を代数方程式へ移す。

\subsection{カタラン数：非線形漸化式を解く}

カタラン数 $C_n$ を、
\[
  C_0=1,\qquad
  C_{n+1}=\sum_{k=0}^{n}C_kC_{n-k}
\]
で定める。
最初の項は、
\[
  1,1,2,5,14,42,\ldots
\]
である。

母関数
\[
  C(x)=\sum_{n=0}^{\infty}C_nx^n
\]
を考える。
漸化式の右辺は畳み込みなので、母関数では積になる。
\[
  C(x)-1=xC(x)^2
\]
すなわち、
\[
  xC(x)^2-C(x)+1=0
\]
である。

二次方程式として解くと、
\[
  C(x)=\frac{1\pm\sqrt{1-4x}}{2x}
\]
となる。
$x=0$ の近くで $C(x)\to C_0=1$ でなければならないので、正しい枝は、
\[
  C(x)=\frac{1-\sqrt{1-4x}}{2x}
\]
である。

もともとの漸化式は、各項がそれ以前の項の積の和で表される非線形なものだった。
それが、母関数を導入すると二次方程式になった。
母関数は線形漸化式専用の道具ではない。
「組合せの分解が畳み込みになる」場面では、非線形な漸化式にも力を発揮する。

\subsection{母関数から見る漸化式の意味}

母関数による変換は、
\[
  \text{無限個の数}
  \longrightarrow
  \text{一つの形式的な対象}
\]
である。
この変換によって、シフトは $x$ の掛け算になり、畳み込みは積になり、漸化式は代数方程式になる。

何を簡単にしたいかによって、どの対象へ移すかが決まる。
差分では変化を簡単にし、行列では連立した状態を簡単にし、母関数では無限個の項の関係を簡単にする。

\section{分数一次変換と離散リカッチ方程式}

\subsection{一次分数型は関数の反復である}

\[
  a_{n+1}=\frac{pa_n+q}{ra_n+s}
\]
を考える。
これは一次分数変換、またはメビウス変換と呼ばれる。
分母があるため一見扱いにくいが、関数の反復として見ると構造が現れる。

不動点は、
\[
  \alpha=\frac{p\alpha+q}{r\alpha+s}
\]
より、
\[
  r\alpha^2+(s-p)\alpha-q=0
\]
を満たす。
従って、不動点は二次方程式の解として求められる。

\subsection{二つの不動点を基準にする}

異なる二つの不動点を $\alpha,\beta$ とする。
次の変数変換を考える。
\[
  b_n=\frac{a_n-\alpha}{a_n-\beta}
\]
これは、$a_n$ が $\alpha$ と $\beta$ のどちらに近いかを、比として記録する量である。

変換
\[
  f(x)=\frac{px+q}{rx+s}
\]
について、固定点の条件を使うと、
\[
  f(x)-\alpha
  =
  \frac{(p-r\alpha)(x-\alpha)}{rx+s}
\]
および
\[
  f(x)-\beta
  =
  \frac{(p-r\beta)(x-\beta)}{rx+s}
\]
となる。
従って、
\[
\begin{aligned}
  \frac{f(x)-\alpha}{f(x)-\beta}
  &=
  \frac{p-r\alpha}{p-r\beta}
  \frac{x-\alpha}{x-\beta}.
\end{aligned}
\]

ここで
\[
  \lambda=\frac{p-r\alpha}{p-r\beta}
\]
と置けば、
\[
  b_{n+1}=\lambda b_n
\]
となる。

分数を含む複雑な漸化式が、変数変換によって等比数列になった。
本編で扱った一次分数型の解法は、このような「二つの基準点からの比を取る」という幾何学的な意味を持っている。

\subsection{連続リカッチ方程式との共通点}

連続的なリカッチ方程式では、既知の特別解を基準にして、
\[
  x=x_0+\frac1u
\]
と置いた。
離散的な一次分数型では、二つの不動点を基準にして、
\[
  b=\frac{x-\alpha}{x-\beta}
\]
と置いた。

変換の形は違うが、発想は同じである。
\[
  \text{非線形で扱いにくい変化}
  \longrightarrow
  \text{基準となる解や点との差}
  \longrightarrow
  \text{線形または等比的な変化}
\]

漸化式と微分方程式を結ぶのは、公式の対応ではない。
「どの量を記録すれば、変化が単純になるか」という問題意識の対応である。

\section{漸近解析：一般項がなくても構造を知る}

\subsection{単調性と有界性}

数列 $a_n$ が
\[
  a_{n+1}\geq a_n
\]
を満たすとき、単調増加という。
また、ある定数 $M$ が存在して
\[
  a_n\leq M
\]
となるとき、上に有界という。

単調増加かつ上に有界な数列は収束する。
これは実数の連続性に関わる基本的な定理である。
証明の考え方は、数列の値が作る集合の上限 $L$ を取ることにある。
もし $a_n$ が $L$ より離れたままなら、$L$ が上限であることに反する。

漸化式の極限を調べるとき、
\[
  \text{単調性}+\text{有界性}
\]
を確認するのは、この定理を使うためである。
不動点方程式だけでは収束を証明できないが、単調性と有界性を加えると、極限の存在を確保できる場合がある。

\subsection{収束先を不動点方程式で絞る}

\[
  a_{n+1}=f(a_n)
\]
で、$a_n\to L$ だとする。
$f$ が連続なら、
\[
  L=f(L)
\]
が必要である。
まず不動点を求め、次に数列が実際にその不動点へ近づくことを示す。

例えば、
\[
  a_{n+1}=\frac{a_n+2}{2}
\]
なら、不動点は $L=2$ である。
$b_n=a_n-2$ と置くと、
\[
  b_{n+1}=\frac12b_n
\]
だから、
\[
  a_n=2+\frac{a_0-2}{2^n}
\]
となり、確かに $2$ へ収束する。

\subsection{縮小写像としての収束証明}

ある区間 $I$ 上で、
\[
  |f(x)-f(y)|\leq q|x-y|
  \qquad(0<q<1)
\]
が成り立つとする。
このような関数を縮小写像という。

反復 $a_{n+1}=f(a_n)$ に対して、
\[
  |a_{n+1}-a_n|
  \leq q|a_n-a_{n-1}|
  \leq q^n|a_1-a_0|
\]
である。
従って、後の項同士の距離は、等比的に小さくなる。

実際、$m>n$ に対して、
\[
\begin{aligned}
  |a_m-a_n|
  &\leq
  |a_m-a_{m-1}|+\cdots+|a_{n+1}-a_n|\\
  &\leq
  (q^{m-1}+\cdots+q^n)|a_1-a_0|\\
  &\leq
  \frac{q^n}{1-q}|a_1-a_0|.
\end{aligned}
\]
$n$ を大きくすれば右辺は $0$ に近づくので、数列は一つの値へ近づく。
その極限を $L$ とすれば、連続性から $L=f(L)$ である。

これは、単に不動点を見つけるだけでなく、近くの点が縮む仕組みから収束を保証する考え方である。
ニュートン法の局所的な収束や、線形漸化式の $|r|<1$ も、この縮みという見方で統一できる。

\subsection{支配的な項を見つける}

フィボナッチ数列の一般項は、
\[
  F_n=\frac1{\sqrt5}\left(\frac{1+\sqrt5}{2}\right)^n
  -\frac1{\sqrt5}\left(\frac{1-\sqrt5}{2}\right)^n
\]
である。
二つの根の絶対値を比べると、
\[
  \left|\frac{1-\sqrt5}{2}\right|<1
\]
だから、後半の項は $n$ が大きくなると消えていく。
従って、
\[
  F_n\sim
  \frac1{\sqrt5}
  \left(\frac{1+\sqrt5}{2}\right)^n
\]
となる。

一般に、線形漸化式では特性根の絶対値が最大のものが、十分大きな $n$ での成長を支配する。
同じ絶対値の根が複数あれば、振動や $n$ の因子が残る。
一般項を完全に書くことと、長期的な成長を理解することは、関連しているが別の仕事である。

\subsection{漸近記号を使う}

\[
  a_n\sim b_n
\]
とは、
\[
  \lim_{n\to\infty}\frac{a_n}{b_n}=1
\]
という意味である。
また、
\[
  a_n=O(b_n)
\]
とは、ある定数 $C$ と十分大きな $n$ に対して、
\[
  |a_n|\leq C|b_n|
\]
となることをいう。

例えば、
\[
  n^2+3n+1\sim n^2
\]
である。
最高次の項が成長を支配するからである。
漸化式の解析では、正確な一般項を得る前に、どの程度の大きさかを見積もることが有効になる。

\section{漸化式が組合せ論と確率へ進む}

\subsection{パスカルの三角形}

二項係数は、
\[
  \binom{n+1}{k}
  =
  \binom{n}{k-1}+\binom{n}{k}
\]
を満たす。
これは、$(n+1)$ 個から $k$ 個を選ぶとき、特別な一個を選ぶ場合と選ばない場合に分けたものである。

二項係数は、二次元の漸化式である。
一方向に並んだ数列だけでなく、格子上の位置 $(n,k)$ を動く規則も漸化式と考えられる。

この漸化式を繰り返すと、
\[
  \binom{n}{0},\binom{n}{1},\ldots,\binom{n}{n}
\]
ができ、パスカルの三角形になる。
さらに、
\[
  \sum_{k=0}^{n}\binom{n}{k}x^k=(1+x)^n
\]
という母関数的な式も現れる。

漸化式、組合せ、二項定理は、別々の章の内容ではない。
「一つの対象を、最後の選択によって場合分けする」という同じ考え方を持っている。

\subsection{確率遷移と行列}

二つの状態 $A,B$ を行き来する確率を考える。
時刻 $n$ に状態 $A,B$ にいる確率を $u_n,v_n$ とし、
\[
  \begin{pmatrix}u_{n+1}\\v_{n+1}\end{pmatrix}
  =
  \begin{pmatrix}p&q\\1-p&1-q\end{pmatrix}
  \begin{pmatrix}u_n\\v_n\end{pmatrix}
\]
とする。
各列または各行の和が $1$ になる行列を、確率の遷移行列として使う。

これは本編で見た連立漸化式と同じ形である。
行列の固有値のうち、$1$ 以外のものの絶対値が $1$ より小さければ、初期状態の影響が消え、定常的な確率分布へ近づく。

数列の極限、行列の固有値、確率過程の定常分布が、同じ線形代数の言葉で説明される。
漸化式は、確率を扱うための最も基本的な時間発展の記法でもある。

\subsection{数論の中の漸化式}

ユークリッドの互除法も、余りを次々に計算する過程として書ける。
\[
  a_{n-1}=q_na_n+a_{n+1},
  \qquad 0\leq a_{n+1}<a_n
\]
である。
余りが小さくなるため、有限回で停止する。

ここでは一般項を求めていない。
それでも、各段階の更新規則と、値が単調に小さくなることから、アルゴリズムの停止を証明できる。

漸化式の考え方は、閉じた公式を得るためだけでなく、計算手順が必ず終わることや、計算量を見積もることにも使われる。

\section{解けない漸化式と、解くことの意味}

\subsection{「解けない」は無知という意味ではない}

本書でいう「解ける」とは、主に高校数学で扱う関数と有限回の操作によって一般項を表せることである。
従って「解けない」とは、数学的に何も分からないという意味ではない。

例えば、
\[
  a_{n+1}=a_n^2+1,\qquad a_0=0
\]
なら、
\[
  0,1,2,5,26,677,\ldots
\]
と順番に計算できる。
一般項が簡単な式で書けなくても、次の問いは残っている。

\begin{itemize}
  \item すべての項は正か。
  \item 単調に増加するか。
  \item どの程度の速さで増えるか。
  \item ある範囲に留まるか。
  \item 初期値を変えると、どのように動きが変わるか。
\end{itemize}

一般項は、漸化式を理解するための一つの答えにすぎない。

\subsection{数学では、問いを変えることが解法になる}

一般項が求めにくいとき、同じ問いに固執する必要はない。
\[
  \text{正確な値を求める}
\]
から、
\[
  \text{増加の速さを求める}
\]
へ変えてもよい。
さらに、
\[
  \text{どの対象へ移せば構造が見えるか}
\]
を問うてもよい。

差分は変化を見えるようにした。
ニュートン法は、解そのものではなく解への道を作った。
母関数は無限個の項を一つにまとめた。
行列は複数の数列を一つの状態にまとめた。
チェビシェフ多項式は、三角関数の振動を多項式へ翻訳した。
漸近解析は、正確な値から長期的な傾向へ焦点を移した。

これらは、ばらばらなテクニックではない。
扱いにくい対象を見つけたとき、何を保存し、何を捨て、どの形へ移せばよいかを考えるという、同じ数学的態度である。

\section{付録の地図を閉じる}

この付録では、漸化式を次のような一つの流れとして眺めた。

\[
\begin{array}{c}
\text{隣り合う項の差} \\
\downarrow \\
\text{差分と離散的な積分} \\
\downarrow \\
\text{反復と力学系} \\
\downarrow \\
\text{微分方程式と数値計算} \\
\downarrow \\
\text{ニュートン法と近似} \\
\downarrow \\
\text{行列と固有値} \\
\downarrow \\
\text{母関数と組合せ} \\
\downarrow \\
\text{三角関数とチェビシェフ多項式} \\
\downarrow \\
\text{不動点と漸近解析} \\
\downarrow \\
\text{確率・数論・アルゴリズム}
\end{array}
\]

最初は、漸化式は数列の項を計算するための式に見える。
しかし、項と項の関係を丁寧に見ると、変化、反復、近似、対称性、線形性、極限という数学の基本概念が姿を現す。

漸化式を解くとは、一般項を一つ書くことだけではない。
その規則が何を記録し、どこへつながり、どの情報を残しているのかを読むことである。
式を見たときに、すぐ計算を始めるのではなく、
\[
  \text{この規則を、別の数学の言葉で言い換えられないか}
\]
と考えてみる。
そこから先に、漸化式の本当の広がりがある。
